\documentclass[conference]{IEEEtran}
\IEEEoverridecommandlockouts

\usepackage[utf8]{inputenc}
\usepackage[T2A]{fontenc}
\usepackage[russian]{babel}

\usepackage{cite}
\usepackage{amsmath,amssymb,amsfonts}
\usepackage{algorithmic}
\usepackage{graphicx}
\usepackage{textcomp}
\usepackage{xcolor}

\def\BibTeX{{\rm B\kern-.05em{\sc i\kern-.025em b}\kern-.08em
    T\kern-.1667em\lower.7ex\hbox{E}\kern-.125emX}}

\newcommand{\R}{\mathbb{R}}
\newcommand{\E}{\mathbb{E}}
\DeclareMathOperator*{\argmin}{arg\,min}

\begin{document}

\title{Beyond Gaussian Noise: Convergence of Stochastic
Optimizers under Correlated and Heavy-Tailed Financial Noise\\
\large (Project Proposal)}

\author{%
\IEEEauthorblockN{Дмитрий\,Кортунов}
\IEEEauthorblockA{\textit{НИУ ВШЭ}\\
Москва, Россия\\
dvkortunov@edu.hse.ru}\\
\and
\IEEEauthorblockN{Вадим\,Пастушенко}
\IEEEauthorblockA{\textit{НИУ ВШЭ}\\
Москва, Россия\\
vapastushenko@edu.hse.ru}
\and
\IEEEauthorblockN{Алексей\,Дубовицкий}
\IEEEauthorblockA{\textit{НИУ ВШЭ}\\
Москва, Россия\\
andubovitskiy@edu.hse.ru}
\and
\IEEEauthorblockN{Андрей\,Шадрин}
\IEEEauthorblockA{\textit{НИУ ВШЭ}\\
Москва, Россия\\
arshadrin@edu.hse.ru}
}

\maketitle

\begin{abstract}
Финансовые ряды содержат шум различной природы --- гауссов,
тяжёлохвостовый, коррелированный. Прямая подача в стандартные
стохастические оптимизаторы нарушает предпосылки теории сходимости. Мы
рассматриваем шумоподавление наблюдаемого ряда как явный этап
оптимизационной процедуры и формулируем проверяемую гипотезу:
\emph{нелинейная порядковая статистика (оконная медиана) восстанавливает
конечную эффективную дисперсию градиентного шума там, где линейные
фильтры этого сделать не могут, и тем самым спасает сходимость SGD при
$\alpha$-устойчивом шуме}. Этот документ --- план работ; полный вывод
механизма, доказательства и численные результаты вынесены в
сопутствующую статью того же формата (\texttt{paper.pdf} /
\texttt{paper\_ru.pdf}), которую настоящий план дополняет.
\end{abstract}

\begin{IEEEkeywords}
шумоподавление, финансовые ряды, стохастическая оптимизация,
$\alpha$-устойчивый шум, медианная фильтрация, фильтр Калмана,
вейвлет-порог, сходимость.
\end{IEEEkeywords}

\section{Введение}

При обучении на финансовых рядах фокусируются обычно на архитектуре
модели или модификации оптимизатора, оставляя структуру входного шума
«за кадром». Между тем именно она --- тяжёлые хвосты, долгая память,
выбросы --- определяет сложность обучения: нарушаются предпосылки
($\E\|\xi\|^2<\infty$, независимость), при которых гарантирована
стабильная работа классических методов \cite{cont2001,mandelbrot1963,
simsekli2019tailindex}. Мы рассматриваем шумоподавление не как
вспомогательную визуализацию, а как полноценную часть процедуры
вычисления градиентного шага, и проверяем, может ли корректная
предобработка изменить свойства градиентного шума так, чтобы метод
сходился там, где без неё расходится. Ключевой тезис, который проект
формулирует и проверяет: разделение фильтров на \emph{линейные}
(скользящее среднее, Калман --- замкнуты относительно $\alpha$-устойчивого
закона и потому бессильны против тяжёлых хвостов) и \emph{нелинейные
порядковые} (медиана --- отображает окно с бесконечной дисперсией в
выход с конечной).

\section{Постановка}

\subsection{Математическая постановка}
Решается задача обучения. Параметры $x\in\R^d$; из наблюдаемого ряда
$y_{1:T}$ строятся пары $(u_i,v_i)_{i=1}^n$, где вход $u_i\in\R^p$ ---
лаговое окно, цель $v_i$ --- следующее значение (или его знак). При
модели $g_x$ и потере $\ell$ (квадратичной для регрессии,
логистической для классификации) минимизируется эмпирический риск
\begin{equation}
f(x)=\frac1n\sum_{i=1}^n\ell\!\big(g_x(u_i),v_i\big),\qquad
x^\star\in\argmin_{x\in\R^d}f(x). \label{eq:risk}
\end{equation}
Оптимизация --- стохастическим методом первого порядка: на шаге $k$
доступен зашумлённый градиент
\begin{equation}
g_k=\nabla f(x_k)+\xi_k,\qquad x_{k+1}=x_k-\eta_k\,\phi(g_k),
\label{eq:step}
\end{equation}
$\xi_k$ --- центрированный шум, без предположений ограниченности
$\E\|\xi_k\|^2$ и независимости. Вход всей процедуры --- наблюдаемый
ряд $y_t=s_t+\xi_t$ ($s$ --- сигнал, $\xi$ --- шум); перед оптимизатором
он пропускается через фильтр $\hat s=F(y_{1:T})$, оптимизатор работает
с $\hat s$; выход --- параметры $x_K$ и метрики. Контроль ---
$F_0=\mathrm{Id}$. Фильтр стоит строго между наблюдением и
стохастическим градиентом~\eqref{eq:step}.

\subsection{Классы шума}
\textbf{N1} гауссов белый $\xi_t\sim\mathcal N(0,\sigma^2)$ (контроль).
\textbf{N2} розовый: $\xi_t=\sum_{j\ge0}\psi_j\eta_{t-j}$, $\psi_j$ ---
FARIMA$(0,d,0)$, $d\in(0,\tfrac12)$ (конечная дисперсия, долгая
память). \textbf{N3} тяжёлохвостовый $\xi_t\sim
\mathcal S_\alpha(\sigma,0,0)$, $\alpha\in(1,2)$ (конечен лишь $p$-й
момент при $p<\alpha$; дисперсия бесконечна). \textbf{N4} смешанный:
$1/f$-фильтр от $\alpha$-устойчивых инноваций (тяжёлые хвосты и долгая
память одновременно). Определяющее свойство: у N3/N4 второй момент
бесконечен, у N1/N2 --- конечен.

\subsection{Фильтры}
Рассматриваются \emph{только} F0--F4. \textbf{F0} тождество.
\textbf{F1} причинное скользящее среднее
$\hat s_t=\tfrac1w\sum_{i=0}^{w-1}y_{t-i}$ (линеен; конечная линейная
комбинация $\alpha$-устойчивых снова $\alpha$-устойчива). \textbf{F2}
Калман (локальный уровень) --- линейная БИХ-рекурсия, оптимальная при
гауссовом шуме \cite{kalman1960}. \textbf{F3} мягкий вейвлет-порог
Добеши $\lambda=\hat\sigma\sqrt{2\ln N}$
\cite{donoho1994wavelets,mallat1999wavelet}. \textbf{F4} причинная
медиана окна $w$ --- единственный нелинейный, порядковая статистика,
устойчива к выбросам \cite{yin1996median}. На реальных данных ---
только причинные F0, F1, F2, F4 (без заглядывания вперёд).

\subsection{Оптимизаторы}
SGD (постоянный шаг $\phi(g)=g$); Clipped-SGD; Normalized-SGD; Adam
\cite{kingma2015adam}; AdamW \cite{loshchilov2019adamw}. Опорный факт:
для SGD с постоянным шагом на $L$-гладкой $\mu$-сильно выпуклой $f$ при
$\E\|\xi_k\|^2\le\sigma^2$ возникает конечное плато
$\E\|\nabla f\|^2\to\Theta(\eta L\sigma^2)$ \cite{robbins1951sgd};
при $\sigma^2=\infty$ плато не определено --- мотив всей работы.

\subsection{Метрики}
Не одна метрика, а батарея (медиана и $95\%$ bootstrap-CI по сидам):
бинарная сходимость (доля сошедшихся прогонов), $T(\varepsilon)=
\min\{K:\E\|\nabla f(x_K)\|^2\le\varepsilon\}$, шумовой пол
(перцентили), наклон $\log\|g\|^2$ vs $\log k$, AUC кривой,
holdout-качество (без смешения train/test), искажение
$\hat H,\hat\alpha$, wall-clock с учётом стоимости фильтра.

\section{Гипотеза и ожидаемый механизм}

Центральная гипотеза: при N3 ($\alpha<2$) выборочная медиана окна
$w$ имеет конечную дисперсию $\tfrac{1}{4wh(0)^2}+o(1/w)$ для
\emph{любого} $\alpha$ (ЦПТ медианы использует лишь $h(0)>0$, не
моменты), тогда как любая линейная комбинация $\alpha$-устойчивых
остаётся $\alpha$-устойчивой с бесконечной дисперсией. Отсюда
предсказания, которые проект должен подтвердить или опровергнуть:
(i)~SGD без фильтра и с F1/F2/F3 не сходится при тяжёлых хвостах;
(ii)~с F4 сходится, шумовой пол падает на $1$--$2$ порядка;
(iii)~на чистом гауссовом N1 медиана \emph{вредит} (потеря
эффективности $\pi/2$); (iv)~на смешанном N4 медиана не спасает
(зависимость окна ломает аргумент). Эти предсказания и их проверка ---
содержание сопутствующей статьи.

\section{Обзор литературы}

Совместный тяжёлохвостовый/долгопамятный шум в стохастической
аппроксимации впервые разобран Anantharam и Borkar
\cite{anantharam2012sa} --- асимптотический анализ через предельное
ОДУ; Chandak и др. \cite{chandak2026hl} дают первые конечновременные
оценки моментов для того же режима. Обе работы анализируют
немодифицированную итерацию. Робастность к тяжёлым хвостам со стороны
оптимизатора: обрезание \cite{gorbunov2020clipping}, high-probability
оценки \cite{cutkosky2021high}. Классические операторы шумоподавления
--- Калман \cite{kalman1960}, вейвлет-порог
\cite{donoho1994wavelets,mallat1999wavelet}, медиана
\cite{yin1996median}; шумоподавление как предобработка прогноза ---
эмпирически \cite{tang2021denoising}; тяжёлый градиентный шум
нейросетей --- \cite{simsekli2019tailindex}. Наш вопрос ортогонален:
при фиксированном оптимизаторе --- \emph{где и насколько} предобработка
наблюдаемого ряда спасает сходимость, с явным механизмом.

\section{План работ}

\textit{Этап~1.} Генераторы N1--N4; проверка, что обратные оценки
$\hat\alpha,\hat H$ восстанавливают заданные параметры.
\textit{Этап~2.} Реализация F0--F4 как причинных отображений;
валидация по SNR и сохранению структуры.
\textit{Этап~3.} Фиксированный набор оптимизаторов; контрольное
сравнение F0 с F1--F4 по батарее метрик.
\textit{Этап~4.} Две серии: (а)~синтетика (квадратичная,
логистическая, AR) --- факторный дизайн (F0--F4)\,$\times$\,(N1--N4)\,
$\times$\,(оптимизатор), $8$ сидов; (б)~реальные ряды нескольких
доменов и частот при той же причинной процедуре (walk-forward, без
lookahead).
\textit{Этап~5.} Калибровка синтетики под измеренные на реальных рядах
$\hat\alpha,\hat H$; статистическое сопоставление; вывод механизма;
отчёт и сопутствующая статья.

\section{Ожидаемые результаты}

Ожидается (и в сопутствующей статье подтверждено с CI по сидам):
матрица сходимости/качества для всех (F0--F4, N1--N4, оптимизатор);
доказанный механизм спасения сходимости медианой и его границы;
практическая таблица «при каком шуме что применять»; воспроизводимый
код. Положительным итогом является связка \emph{(причинная медиана~+
SGD)} как процедура, делающая обучение при тяжёлохвостовом шуме
сходящимся там, где без предобработки оно не сходится; негативные
зоны (чистый гаусс, смешанный долгопамятный режим, сырые near-white
доходности) сообщаются честно.

\begin{thebibliography}{00}

\bibitem{kalman1960}
R.~E. Kalman, ``A new approach to linear filtering and prediction
problems,'' \textit{J. Basic Eng.}, vol.~82, no.~1, pp.~35--45, 1960.

\bibitem{donoho1994wavelets}
D.~L. Donoho and I.~M. Johnstone, ``Ideal spatial adaptation by wavelet
shrinkage,'' \textit{Biometrika}, vol.~81, no.~3, pp.~425--455, 1994.

\bibitem{mallat1999wavelet}
S.~Mallat, \textit{A Wavelet Tour of Signal Processing}, 2nd~ed. San
Diego: Academic Press, 1999.

\bibitem{yin1996median}
L.~Yin, R.~Yang, M.~Gabbouj, and Y.~Neuvo, ``Weighted median filters: a
tutorial,'' \textit{IEEE Trans. Circuits Syst. II}, vol.~43, no.~3,
pp.~157--192, 1996.

\bibitem{cont2001}
R.~Cont, ``Empirical properties of asset returns: stylized facts and
statistical issues,'' \textit{Quant. Finance}, vol.~1, no.~2,
pp.~223--236, 2001.

\bibitem{mandelbrot1963}
B.~B. Mandelbrot, ``The variation of certain speculative prices,''
\textit{J. Business}, vol.~36, no.~4, pp.~394--419, 1963.

\bibitem{simsekli2019tailindex}
U.~\c{S}im\c{s}ekli, L.~Sagun, and M.~G\"{u}rb\"{u}zbalaban, ``A
tail-index analysis of stochastic gradient noise in deep neural
networks,'' in \textit{Proc. ICML}, 2019, pp.~5827--5837.

\bibitem{robbins1951sgd}
H.~Robbins and S.~Monro, ``A stochastic approximation method,''
\textit{Ann. Math. Stat.}, vol.~22, no.~3, pp.~400--407, 1951.

\bibitem{kingma2015adam}
D.~P. Kingma and J.~Ba, ``Adam: a method for stochastic
optimization,'' in \textit{Proc. ICLR}, 2015.

\bibitem{loshchilov2019adamw}
I.~Loshchilov and F.~Hutter, ``Decoupled weight decay
regularization,'' in \textit{Proc. ICLR}, 2019.

\bibitem{gorbunov2020clipping}
E.~Gorbunov, M.~Danilova, and A.~Gasnikov, ``Stochastic optimization
with heavy-tailed noise via accelerated gradient clipping,'' in
\textit{Proc. NeurIPS}, vol.~33, 2020, pp.~15042--15053.

\bibitem{cutkosky2021high}
A.~Cutkosky and H.~Mehta, ``High-probability bounds for non-convex
stochastic optimization with heavy tails,'' in \textit{Proc. NeurIPS},
2021.

\bibitem{anantharam2012sa}
V.~Anantharam and V.~S. Borkar, ``Stochastic approximation with long
range dependent and heavy tailed noise,'' \textit{Queueing Syst.},
vol.~71, no.~1--2, pp.~221--242, 2012,
doi:~10.1007/s11134-012-9283-0.

\bibitem{chandak2026hl}
S.~Chandak, A.~K. Yadav, A.~\"{O}zg\"{u}r, and N.~Bambos,
``Heavy-tailed and long-range dependent noise in stochastic
approximation: a finite-time analysis,'' arXiv:2603.19648, Mar.\ 2026.

\bibitem{tang2021denoising}
Q.~Tang, T.~Fan, R.~Shi, J.~Huang, and Y.~Ma, ``Prediction of
financial time series using LSTM and data denoising methods,''
arXiv:2103.03505, 2021.

\end{thebibliography}

\end{document}
